%% Generated by tools/gen_error_codes.py from docs/errors/ of the compiler;
%% edit the compiler's pages or tools/error_themes.txt, then rerun it.

\clearpage
\section{Laziness, generators and comprehensions}
\label{sec:codes:lazy}

The book's chapter \textit{Laziness} specifies the rules behind these codes.

\begin{longtable}{@{}l p{0.78\linewidth} r@{}}
\toprule Code & Message & Page\\ \midrule \endhead
\hyperref[sec:codes:e4087]{E4087} & cannot construct a lazy value of type \errhole{} implicitly & \pageref{sec:codes:e4087}\\
\hyperref[sec:codes:e4108]{E4108} & validation of lazy closure function fails to generate value of type \errhole{} & \pageref{sec:codes:e4108}\\
\hyperref[sec:codes:e4109]{E4109} & the size of the list comprehension, constructed from type \errhole{}, must be determinable at compile time & \pageref{sec:codes:e4109}\\
\hyperref[sec:codes:e4249]{E4249} & the construction of a lazy value has no effect & \pageref{sec:codes:e4249}\\
\hyperref[sec:codes:e4291]{E4291} & yield statement is not within a generator & \pageref{sec:codes:e4291}\\
\hyperref[sec:codes:e4292]{E4292} & a yield statement must yield a value of type \errhole{} & \pageref{sec:codes:e4292}\\
\hyperref[sec:codes:e4293]{E4293} & cannot yield inside a try block, a catch or a scope guard & \pageref{sec:codes:e4293}\\
\hyperref[sec:codes:e4294]{E4294} & a generator cannot return a value, its values are yielded & \pageref{sec:codes:e4294}\\
\hyperref[sec:codes:e4295]{E4295} & a generator cannot throw exceptions & \pageref{sec:codes:e4295}\\
\hyperref[sec:codes:e4296]{E4296} & the parameter \errhole{} of a generator cannot be declared \errhole{} & \pageref{sec:codes:e4296}\\
\hyperref[sec:codes:e4297]{E4297} & a generator cannot hold a value of type \errhole{}, which has a destructor & \pageref{sec:codes:e4297}\\
\hyperref[sec:codes:e4299]{E4299} & yield in expects a generator, not a value of type \errhole{} & \pageref{sec:codes:e4299}\\
\hyperref[sec:codes:e4300]{E4300} & copying a generator comprehension has no effect & \pageref{sec:codes:e4300}\\
\hyperref[sec:codes:e4301]{E4301} & cannot \errhole{} out of the element of a generator comprehension & \pageref{sec:codes:e4301}\\
\hyperref[sec:codes:e4302]{E4302} & the size of a list comprehension filtered by if is not determinable at compile time & \pageref{sec:codes:e4302}\\
\hyperref[sec:codes:e4303]{E4303} & the filter of the list comprehension, constructed from type \errhole{}, must be determinable at compile time & \pageref{sec:codes:e4303}\\
\bottomrule
\end{longtable}

\errorcode{E4087}{cannot construct a lazy value of type \errhole{} implicitly}
\label{sec:codes:e4087}

Raised during \textbf{validation}, from \texttt{ValidateErrorMessage::\allowbreak{}IMPLICIT\_\allowbreak{}LAZY} (\textit{src/\allowbreak{}ymirc/\allowbreak{}semantic/\allowbreak{}validator/\allowbreak{}errors.yr}).

\subsubsection*{What it means}

A value was used where a lazy one is expected, without the \tokennolst{lazy} keyword. A lazy value defers its computation, which changes when the expression runs, so it is never built implicitly.

\subsubsection*{Example}

\textit{test\_\allowbreak{}resources/\allowbreak{}diagnostics/\allowbreak{}test41.yr}:

%% check: skip
\begin{lstlisting}[style=coloredverbatimError]
fn foo (lazy x: i32)-> i32 { x }

fn main () { foo (1); }
\end{lstlisting}

\begin{lstlisting}[style=bashVerb, breaklines=true, escapechar=~]
Error[E4087] : cannot construct a lazy value of type i32 implicitly
 --> test_resources/diagnostics/test41.yr:(3,19)
 3  ~\lstbox{\textSFxi{}}~ fn main () { foo (1); }
    ~\lstbox{\textSFv{}}~                   ^  note: for parameter lazy x: i32
    ~\lstbox{\textSFii{}}~~\lstbox{\textSFx{}}~~\lstbox{\textSFx{}}~~\lstbox{\textSFx{}}~~\lstbox{\textSFx{}}~~\lstbox{\textSFx{}}~~\lstbox{\textSFx{}}~~\lstbox{\textSFx{}}~
\end{lstlisting}

\subsubsection*{How to fix it}

Write \tokennolst{lazy} before the expression.

\errorcode{E4108}{validation of lazy closure function fails to generate value of type \errhole{}}
\label{sec:codes:e4108}

Raised during \textbf{validation}, from \texttt{ValidateErrorMessage::\allowbreak{}LAZY\_\allowbreak{}VALIDATION} (\textit{src/\allowbreak{}ymirc/\allowbreak{}semantic/\allowbreak{}validator/\allowbreak{}errors.yr}).

\subsubsection*{What it means}

A lazy value could not be built: validating the closure that defers the computation failed to produce a value of the expected type.

\subsubsection*{Example}

\textit{test\_\allowbreak{}resources/\allowbreak{}diagnostics/\allowbreak{}test109.yr}:

%% check: skip
\begin{lstlisting}[style=coloredverbatimError]
// a lazy value assigning a variable it encloses
fn main () {
    let mut x = 1;
    let y = lazy { x = 2; 1 };
    let _ = y;
    x = 3;
}
\end{lstlisting}

\begin{lstlisting}[style=bashVerb, breaklines=true, escapechar=~]
Error[E4108] : validation of lazy closure function fails to generate value of type mut i32
 --> test_resources/diagnostics/test109.yr:(4,13)
 4  ~\lstbox{\textSFxi{}}~     let y = lazy { x = 2; 1 };
    ~\lstbox{\textSFv{}}~             ^^^^   -  error: x is enclosed, and cannot be used as a lvalue
    ~\lstbox{\textSFxi{}}~
    = when validating test109::main -- (test_resources/diagnostics/test109.yr:(2,4))
    ~\lstbox{\textSFii{}}~~\lstbox{\textSFx{}}~~\lstbox{\textSFx{}}~~\lstbox{\textSFx{}}~~\lstbox{\textSFx{}}~~\lstbox{\textSFx{}}~~\lstbox{\textSFx{}}~~\lstbox{\textSFx{}}~
Error[E4095] : incompatible types error and error
 --> test_resources/diagnostics/test109.yr:(5,9)
 5  ~\lstbox{\textSFxi{}}~     let _ = y;
    ~\lstbox{\textSFv{}}~         ^   -
    ~\lstbox{\textSFxi{}}~
    = when validating test109::main -- (test_resources/diagnostics/test109.yr:(2,4))
    ~\lstbox{\textSFii{}}~~\lstbox{\textSFx{}}~~\lstbox{\textSFx{}}~~\lstbox{\textSFx{}}~~\lstbox{\textSFx{}}~~\lstbox{\textSFx{}}~~\lstbox{\textSFx{}}~~\lstbox{\textSFx{}}~
\end{lstlisting}

\subsubsection*{How to fix it}

Check the expression the \tokennolst{lazy} applies to; the errors under this diagnostic say what went wrong inside it.

\errorcode{E4109}{the size of the list comprehension, constructed from type \errhole{}, must be determinable at compile time}
\label{sec:codes:e4109}

Raised during \textbf{validation}, from \texttt{ValidateErrorMessage::\allowbreak{}LIST\_\allowbreak{}COMPR\_\allowbreak{}SIZE\_\allowbreak{}CTE} (\textit{src/\allowbreak{}ymirc/\allowbreak{}semantic/\allowbreak{}validator/\allowbreak{}errors.yr}).

\subsubsection*{What it means}

The length of a list comprehension is not computable at compile time, and the type it is being built into needs a statically known size.

\subsubsection*{Example}

\textit{test\_\allowbreak{}resources/\allowbreak{}for\_\allowbreak{}loops/\allowbreak{}lst\_\allowbreak{}compr/\allowbreak{}test12.yr}:

%% check: skip
\begin{lstlisting}[style=coloredverbatimError]
fn main () {
    let mut a = copy [1, 2, 3];
    let _ = [i for i in a];
}
\end{lstlisting}

\begin{lstlisting}[style=bashVerb, breaklines=true, escapechar=~]
Error[E4109] : the size of the list comprehension, constructed from type mut [i32], must be determinable at compile time
 --> test_resources/for_loops/lst_compr/test12.yr:(3,25)
 3  ~\lstbox{\textSFxi{}}~     let _ = [i for i in a];
    ~\lstbox{\textSFv{}}~                         ^  note: did you mean to enclose it within a copy?
    ~\lstbox{\textSFii{}}~~\lstbox{\textSFx{}}~~\lstbox{\textSFx{}}~~\lstbox{\textSFx{}}~~\lstbox{\textSFx{}}~~\lstbox{\textSFx{}}~~\lstbox{\textSFx{}}~~\lstbox{\textSFx{}}~
\end{lstlisting}

\subsubsection*{How to fix it}

Build a slice, whose length is a run time value, or make the bounds of the comprehension compile time constants.

\errorcode{E4249}{the construction of a lazy value has no effect}
\label{sec:codes:e4249}

Raised during \textbf{validation}, from \texttt{ValidateErrorMessage::\allowbreak{}UNECESSARY\_\allowbreak{}LAZY} (\textit{src/\allowbreak{}ymirc/\allowbreak{}semantic/\allowbreak{}validator/\allowbreak{}errors.yr}).

\subsubsection*{What it means}

\tokennolst{lazy} was applied where the value is used immediately, so nothing is deferred.

\subsubsection*{Example}

\textit{test\_\allowbreak{}resources/\allowbreak{}lazyness/\allowbreak{}test16.yr}:

%% check: skip
\begin{lstlisting}[style=coloredverbatimError]
fn main () {
    match (lazy 12) {
        x => { x; }
    }
}
\end{lstlisting}

\begin{lstlisting}[style=bashVerb, breaklines=true, escapechar=~]
Error[E4249] : the construction of a lazy value has no effect
 --> test_resources/lazyness/test16.yr:(2,12)
 2  ~\lstbox{\textSFxi{}}~     match (lazy 12) {
    ~\lstbox{\textSFv{}}~            ^^^^
    ~\lstbox{\textSFii{}}~~\lstbox{\textSFx{}}~~\lstbox{\textSFx{}}~~\lstbox{\textSFx{}}~~\lstbox{\textSFx{}}~~\lstbox{\textSFx{}}~~\lstbox{\textSFx{}}~~\lstbox{\textSFx{}}~
\end{lstlisting}

\subsubsection*{How to fix it}

Drop the \tokennolst{lazy}.

\errorcode{E4291}{yield statement is not within a generator}
\label{sec:codes:e4291}

Raised during \textbf{validation}, from \texttt{ValidateErrorMessage::\allowbreak{}YIELD\_\allowbreak{}NO\_\allowbreak{}GENERATOR} (\textit{src/\allowbreak{}ymirc/\allowbreak{}semantic/\allowbreak{}validator/\allowbreak{}errors.yr}).

\subsubsection*{What it means}

A \tokennolst{yield} appears in a function that is not a generator. Only a function introduced by \tokennolst{yield} can yield values; a lambda nested in a generator is a function of its own, and cannot yield on behalf of the generator either.

\subsubsection*{Example}

\textit{test\_\allowbreak{}resources/\allowbreak{}generators/\allowbreak{}test6.yr}:

%% check: skip
\begin{lstlisting}[style=coloredverbatimError]
// yield outside a generator
fn notGenerator ()-> i32 {
    yield 1;
    2
}
\end{lstlisting}

\begin{lstlisting}[style=bashVerb, breaklines=true, escapechar=~]
Error[E4291] : yield statement is not within a generator
 --> test_resources/generators/test6.yr:(3,5)
 3  ~\lstbox{\textSFxi{}}~     yield 1;
    ~\lstbox{\textSFv{}}~     ^^^^^
    ~\lstbox{\textSFxi{}}~
    = when validating test6::notGenerator -- (test_resources/generators/test6.yr:(2,4))
    ~\lstbox{\textSFii{}}~~\lstbox{\textSFx{}}~~\lstbox{\textSFx{}}~~\lstbox{\textSFx{}}~~\lstbox{\textSFx{}}~~\lstbox{\textSFx{}}~~\lstbox{\textSFx{}}~~\lstbox{\textSFx{}}~
\end{lstlisting}

\subsubsection*{How to fix it}

Introduce the function with \tokennolst{yield} instead of \tokennolst{fn}, its declared return type being the type of the yielded values.

\errorcode{E4292}{a yield statement must yield a value of type \errhole{}}
\label{sec:codes:e4292}

Raised during \textbf{validation}, from \texttt{ValidateErrorMessage::\allowbreak{}YIELD\_\allowbreak{}NO\_\allowbreak{}VALUE} (\textit{src/\allowbreak{}ymirc/\allowbreak{}semantic/\allowbreak{}validator/\allowbreak{}errors.yr}).

\subsubsection*{What it means}

A \tokennolst{yield} without a value hands nothing to the consumer of the generator. The type in the message is the one the generator declares.

\subsubsection*{Example}

\textit{test\_\allowbreak{}resources/\allowbreak{}generators/\allowbreak{}test7.yr}:

%% check: skip
\begin{lstlisting}[style=coloredverbatimError]
// a yield without a value
yield gen ()-> i32 {
    yield;
}
\end{lstlisting}

\begin{lstlisting}[style=bashVerb, breaklines=true, escapechar=~]
Error[E4292] : a yield statement must yield a value of type i32
 --> test_resources/generators/test7.yr:(3,5)
 3  ~\lstbox{\textSFxi{}}~     yield;
    ~\lstbox{\textSFv{}}~     ^^^^^
    ~\lstbox{\textSFxi{}}~
    = when validating test7::gen -- (test_resources/generators/test7.yr:(2,7))
    ~\lstbox{\textSFii{}}~~\lstbox{\textSFx{}}~~\lstbox{\textSFx{}}~~\lstbox{\textSFx{}}~~\lstbox{\textSFx{}}~~\lstbox{\textSFx{}}~~\lstbox{\textSFx{}}~~\lstbox{\textSFx{}}~
\end{lstlisting}

\subsubsection*{How to fix it}

Give the \tokennolst{yield} a value. To end the generator early, use \tokennolst{return;}.

\errorcode{E4293}{cannot yield inside a try block, a catch or a scope guard}
\label{sec:codes:e4293}

Raised during \textbf{validation}, from \texttt{ValidateErrorMessage::\allowbreak{}YIELD\_\allowbreak{}PROTECTED\_\allowbreak{}REGION} (\textit{src/\allowbreak{}ymirc/\allowbreak{}semantic/\allowbreak{}validator/\allowbreak{}errors.yr}).

\subsubsection*{What it means}

A generator is suspended at each \tokennolst{yield} and resumed by a jump back to it. A jump cannot re-enter a try block, a catch or a block guarded by a scope guard, as the handler of the region would not be installed again.

\subsubsection*{Example}

\textit{test\_\allowbreak{}resources/\allowbreak{}generators/\allowbreak{}test8.yr}:

%% check: skip
\begin{lstlisting}[style=coloredverbatimError]
// yield inside a try block
class Err over Exception {
    pub self () {}
}

fn mayThrow ()
    throws Err
{
    throw copy Err ();
}

yield gen ()-> i32 {
    {
        mayThrow ();
        yield 1;
    } catch {
        _: &Err => {}
    }
}
\end{lstlisting}

\begin{lstlisting}[style=bashVerb, breaklines=true, escapechar=~]
Error[E4293] : cannot yield inside a try block, a catch or a scope guard
 --> test_resources/generators/test8.yr:(13,5)
13  ~\lstbox{\textSFxi{}}~     {
    ~\lstbox{\textSFv{}}~     ^
14  ~\lstbox{\textSFxi{}}~         mayThrow ();
15  ~\lstbox{\textSFxi{}}~         yield 1;
    ~\lstbox{\textSFv{}}~         -----
    ~\lstbox{\textSFxi{}}~
    = when validating test8::gen -- (test_resources/generators/test8.yr:(12,7))
    ~\lstbox{\textSFii{}}~~\lstbox{\textSFx{}}~~\lstbox{\textSFx{}}~~\lstbox{\textSFx{}}~~\lstbox{\textSFx{}}~~\lstbox{\textSFx{}}~~\lstbox{\textSFx{}}~~\lstbox{\textSFx{}}~
\end{lstlisting}

\subsubsection*{How to fix it}

Move the \tokennolst{yield} out of the protected region, storing what it yields in a variable of the generator.

\errorcode{E4294}{a generator cannot return a value, its values are yielded}
\label{sec:codes:e4294}

Raised during \textbf{validation}, from \texttt{ValidateErrorMessage::\allowbreak{}RETURN\_\allowbreak{}VALUE\_\allowbreak{}GENERATOR} (\textit{src/\allowbreak{}ymirc/\allowbreak{}semantic/\allowbreak{}validator/\allowbreak{}errors.yr}).

\subsubsection*{What it means}

The declared return type of a generator is the type of the values it yields, and a call of a generator returns its handle. A \tokennolst{return} can only end the generator.

\subsubsection*{Example}

\textit{test\_\allowbreak{}resources/\allowbreak{}generators/\allowbreak{}test10.yr}:

%% check: skip
\begin{lstlisting}[style=coloredverbatimError]
// a generator returning a value
yield gen ()-> i32 {
    yield 1;
    return 2;
}
\end{lstlisting}

\begin{lstlisting}[style=bashVerb, breaklines=true, escapechar=~]
Error[E4294] : a generator cannot return a value, its values are yielded
 --> test_resources/generators/test10.yr:(4,5)
 4  ~\lstbox{\textSFxi{}}~     return 2;
    ~\lstbox{\textSFv{}}~     ^^^^^^
    ~\lstbox{\textSFxi{}}~
    = when validating test10::gen -- (test_resources/generators/test10.yr:(2,7))
    ~\lstbox{\textSFii{}}~~\lstbox{\textSFx{}}~~\lstbox{\textSFx{}}~~\lstbox{\textSFx{}}~~\lstbox{\textSFx{}}~~\lstbox{\textSFx{}}~~\lstbox{\textSFx{}}~~\lstbox{\textSFx{}}~
\end{lstlisting}

\subsubsection*{How to fix it}

Yield the value, then \tokennolst{return;} if the generator has to stop there.

\errorcode{E4295}{a generator cannot throw exceptions}
\label{sec:codes:e4295}

Raised during \textbf{validation}, from \texttt{ValidateErrorMessage::\allowbreak{}GENERATOR\_\allowbreak{}THROWS} (\textit{src/\allowbreak{}ymirc/\allowbreak{}semantic/\allowbreak{}validator/\allowbreak{}errors.yr}).

\subsubsection*{What it means}

The body of a generator runs while its consumer iterates over it, where the exceptions it declares could not be tracked. A generator method, declared by \tokennolst{yield} in a class, a record or a trait, is rejected the same way.

\subsubsection*{Example}

\textit{test\_\allowbreak{}resources/\allowbreak{}generators/\allowbreak{}test11.yr}:

%% check: skip
\begin{lstlisting}[style=coloredverbatimError]
// a generator declaring exceptions
class Err over Exception {
    pub self () {}
}

yield gen ()-> i32
    throws Err
{
    yield 1;
    throw copy Err ();
}
\end{lstlisting}

\begin{lstlisting}[style=bashVerb, breaklines=true, escapechar=~]
Error[E4295] : a generator cannot throw exceptions
 --> test_resources/generators/test11.yr:(7,5)
 7  ~\lstbox{\textSFxi{}}~     throws Err
    ~\lstbox{\textSFv{}}~     ^^^^^^
    ~\lstbox{\textSFxi{}}~
    = when validating test11::gen -- (test_resources/generators/test11.yr:(6,7))
    ~\lstbox{\textSFii{}}~~\lstbox{\textSFx{}}~~\lstbox{\textSFx{}}~~\lstbox{\textSFx{}}~~\lstbox{\textSFx{}}~~\lstbox{\textSFx{}}~~\lstbox{\textSFx{}}~~\lstbox{\textSFx{}}~
\end{lstlisting}

\subsubsection*{How to fix it}

Catch the exceptions inside the generator, outside of the region of any \tokennolst{yield}, cf. \hyperref[sec:codes:e4293]{E4293}.

\errorcode{E4296}{the parameter \errhole{} of a generator cannot be declared \errhole{}}
\label{sec:codes:e4296}

Raised during \textbf{validation}, from \texttt{ValidateErrorMessage::\allowbreak{}GENERATOR\_\allowbreak{}PARAM} (\textit{src/\allowbreak{}ymirc/\allowbreak{}semantic/\allowbreak{}validator/\allowbreak{}errors.yr}).

\subsubsection*{What it means}

The parameters of a generator are stored in its frame, which outlives the call creating it. A \tokennolst{ref} or \tokennolst{lazy} parameter refers to the frame of the caller, that may be gone when the generator is resumed. This applies to generator methods as well; the \tokennolst{self} of a record generator is copied into its frame, so a \tokennolst{mut self}, whose changes the caller would never see, is rejected.

\subsubsection*{Example}

\textit{test\_\allowbreak{}resources/\allowbreak{}generators/\allowbreak{}test12.yr}:

%% check: skip
\begin{lstlisting}[style=coloredverbatimError]
// a generator taking a reference parameter
yield gen (ref x: i32)-> i32 {
    yield x;
}
\end{lstlisting}

\begin{lstlisting}[style=bashVerb, breaklines=true, escapechar=~]
Error[E4296] : the parameter x of a generator cannot be declared ref
 --> test_resources/generators/test12.yr:(2,16)
 2  ~\lstbox{\textSFxi{}}~ yield gen (ref x: i32)-> i32 {
    ~\lstbox{\textSFv{}}~                ^
    ~\lstbox{\textSFxi{}}~
    = when validating test12::gen -- (test_resources/generators/test12.yr:(2,7))
    ~\lstbox{\textSFii{}}~~\lstbox{\textSFx{}}~~\lstbox{\textSFx{}}~~\lstbox{\textSFx{}}~~\lstbox{\textSFx{}}~~\lstbox{\textSFx{}}~~\lstbox{\textSFx{}}~~\lstbox{\textSFx{}}~
\end{lstlisting}

\subsubsection*{How to fix it}

Take the parameter by value.

\errorcode{E4297}{a generator cannot hold a value of type \errhole{}, which has a destructor}
\label{sec:codes:e4297}

Raised during \textbf{validation}, from \texttt{ValidateErrorMessage::\allowbreak{}GENERATOR\_\allowbreak{}MOVABLE} (\textit{src/\allowbreak{}ymirc/\allowbreak{}semantic/\allowbreak{}validator/\allowbreak{}errors.yr}).

\subsubsection*{What it means}

A generator can be abandoned while suspended, and its frame is then collected without its pending destructors ever running. A value whose type has a destructor cannot be held by a generator for that reason. This includes the \tokennolst{self} of a generator method, so an entity with a destructor cannot take the default generator method of a trait.

\subsubsection*{Example}

\textit{test\_\allowbreak{}resources/\allowbreak{}generators/\allowbreak{}test15.yr}:

%% check: skip
\begin{lstlisting}[style=coloredverbatimError]
// a generator holding a value with a destructor
entity Handle {
    pub self () {}

    __dtor (mut self) {}
}

yield gen ()-> i32 {
    let h = Handle ();
    yield 1;
}
\end{lstlisting}

\begin{lstlisting}[style=bashVerb, breaklines=true, escapechar=~]
Error[E4297] : a generator cannot hold a value of type mut test15::Handle, which has a destructor
 --> test_resources/generators/test15.yr:(9,20)
 9  ~\lstbox{\textSFxi{}}~     let h = Handle ();
    ~\lstbox{\textSFv{}}~                    ^
    ~\lstbox{\textSFxi{}}~
    = when validating test15::gen -- (test_resources/generators/test15.yr:(8,7))
    ~\lstbox{\textSFii{}}~~\lstbox{\textSFx{}}~~\lstbox{\textSFx{}}~~\lstbox{\textSFx{}}~~\lstbox{\textSFx{}}~~\lstbox{\textSFx{}}~~\lstbox{\textSFx{}}~~\lstbox{\textSFx{}}~
\end{lstlisting}

\subsubsection*{How to fix it}

Hold the value in the consumer of the generator, and pass the generator what it needs from it.

\errorcode{E4299}{yield in expects a generator, not a value of type \errhole{}}
\label{sec:codes:e4299}

Raised during \textbf{validation}, from \texttt{ValidateErrorMessage::\allowbreak{}YIELD\_\allowbreak{}IN\_\allowbreak{}NO\_\allowbreak{}GENERATOR} (\textit{src/\allowbreak{}ymirc/\allowbreak{}semantic/\allowbreak{}validator/\allowbreak{}errors.yr}).

\subsubsection*{What it means}

\tokennolst{yield in g} yields every value of the generator \tokennolst{g}, until it is over. The value after \tokennolst{in} must therefore be a generator handle, of type \tokennolst{yield-\textgreater{} T}; any other value, even an iterable one, has no values to hand over this way.

\subsubsection*{Example}

\textit{test\_\allowbreak{}resources/\allowbreak{}generators/\allowbreak{}test43.yr}:

%% check: skip
\begin{lstlisting}[style=coloredverbatimError]
// yield in only yields the values of a generator
yield all (a: [i32])-> i32 {
    yield in a;
}

fn main () {
    for x in all (copy [1, 2]) {
        x;
    }
}
\end{lstlisting}

\begin{lstlisting}[style=bashVerb, breaklines=true, escapechar=~]
Error[E4299] : yield in expects a generator, not a value of type [i32]
 --> test_resources/generators/test43.yr:(3,14)
 3  ~\lstbox{\textSFxi{}}~     yield in a;
    ~\lstbox{\textSFv{}}~              ^
    ~\lstbox{\textSFxi{}}~
    = when validating test43::all -- (test_resources/generators/test43.yr:(2,7))
    ~\lstbox{\textSFii{}}~~\lstbox{\textSFx{}}~~\lstbox{\textSFx{}}~~\lstbox{\textSFx{}}~~\lstbox{\textSFx{}}~~\lstbox{\textSFx{}}~~\lstbox{\textSFx{}}~~\lstbox{\textSFx{}}~
\end{lstlisting}

\subsubsection*{How to fix it}

Yield the values one at a time with a loop, \tokennolst{for x in a \{ yield x; \}}, or pass a generator producing them to \tokennolst{yield in}.

\errorcode{E4300}{copying a generator comprehension has no effect}
\label{sec:codes:e4300}

Raised during \textbf{validation}, from \texttt{ValidateErrorMessage::\allowbreak{}UNECESSARY\_\allowbreak{}GENERATOR\_\allowbreak{}COPY} (\textit{src/\allowbreak{}ymirc/\allowbreak{}semantic/\allowbreak{}validator/\allowbreak{}errors.yr}).

\subsubsection*{What it means}

A generator comprehension, \tokennolst{(yield x for x in it)}, builds no list: it creates a generator, which computes each value when it is iterated. There is nothing for \tokennolst{copy} to duplicate, unlike a list comprehension, where \tokennolst{copy} builds a slice instead of an array.

Only a \tokennolst{copy} applied to the comprehension itself is reported. A \tokennolst{dcopy} of a value merely holding one, \texttt{dcopy [(yield x for x in it) for \_ in 0 .. 2]}, does copy the slice, and simply leaves its generators alone.

\subsubsection*{Example}

\textit{test\_\allowbreak{}resources/\allowbreak{}generators/\allowbreak{}test57.yr}:

%% check: skip
\begin{lstlisting}[style=coloredverbatimError]
// a generator comprehension allocates its own frame, there is nothing to copy
fn main () {
    let _ = copy (yield x for x in 0 .. 3);
}
\end{lstlisting}

\begin{lstlisting}[style=bashVerb, breaklines=true, escapechar=~]
Error[E4300] : copying a generator comprehension has no effect
 --> test_resources/generators/test57.yr:(3,18)
 3  ~\lstbox{\textSFxi{}}~     let _ = copy (yield x for x in 0 .. 3);
    ~\lstbox{\textSFv{}}~                  ^
    ~\lstbox{\textSFxi{}}~
    = when validating test57::main -- (test_resources/generators/test57.yr:(2,4))
    ~\lstbox{\textSFii{}}~~\lstbox{\textSFx{}}~~\lstbox{\textSFx{}}~~\lstbox{\textSFx{}}~~\lstbox{\textSFx{}}~~\lstbox{\textSFx{}}~~\lstbox{\textSFx{}}~~\lstbox{\textSFx{}}~
\end{lstlisting}

\subsubsection*{How to fix it}

Drop the \tokennolst{copy}. To collect the values into a slice, use a list comprehension over the generator, \tokennolst{copy [x for x in (yield x for x in it)]}.

\errorcode{E4301}{cannot \errhole{} out of the element of a generator comprehension}
\label{sec:codes:e4301}

Raised during \textbf{validation}, from \texttt{ValidateErrorMessage::\allowbreak{}GENERATOR\_\allowbreak{}COMPR\_\allowbreak{}EXIT} (\textit{src/\allowbreak{}ymirc/\allowbreak{}semantic/\allowbreak{}validator/\allowbreak{}errors.yr}).

\subsubsection*{What it means}

The element of a generator comprehension, \tokennolst{(yield e for x in it)}, is not computed where the comprehension is written, but each time the generator is iterated, possibly long after the enclosing loop or function is over. A \tokennolst{break}, a \tokennolst{continue} or a \tokennolst{return} in it has nothing to leave: it cannot exit the loop or the function around the comprehension, and it cannot skip or end the values of the generator either.

A \tokennolst{yield}, or a \tokennolst{yield in}, is rejected for the same reason read the other way round: the comprehension yields the element and nothing else, so a \tokennolst{yield} written in it would add values to a generator that is not the element's to produce.

A \tokennolst{break} or a \tokennolst{continue} of a loop written inside the element itself is fine.

\subsubsection*{Example, of a \tokennolst{break}}

\textit{test\_\allowbreak{}resources/\allowbreak{}generators/\allowbreak{}test64.yr}:

%% check: skip
\begin{lstlisting}[style=coloredverbatimError]
// the element of a generator comprehension cannot break the loop enclosing the comprehension
fn main (args: [[c8]]) {
    loop {
        let _ = (yield { if (args.len > 2us) break; x } for x in 0 .. 12);
    }
}
\end{lstlisting}

\begin{lstlisting}[style=bashVerb, breaklines=true, escapechar=~]
Error[E4301] : cannot break out of the element of a generator comprehension
 --> test_resources/generators/test64.yr:(4,46)
 4  ~\lstbox{\textSFxi{}}~         let _ = (yield { if (args.len > 2us) break; x } for x in 0 .. 12);
    ~\lstbox{\textSFv{}}~                                              ^^^^^
    ~\lstbox{\textSFxi{}}~
    = when validating test64::main -- (test_resources/generators/test64.yr:(2,4))
    ~\lstbox{\textSFii{}}~~\lstbox{\textSFx{}}~~\lstbox{\textSFx{}}~~\lstbox{\textSFx{}}~~\lstbox{\textSFx{}}~~\lstbox{\textSFx{}}~~\lstbox{\textSFx{}}~~\lstbox{\textSFx{}}~
\end{lstlisting}

\subsubsection*{Example, of a \tokennolst{yield}}

\textit{test\_\allowbreak{}resources/\allowbreak{}generators/\allowbreak{}test68.yr}:

%% check: skip
\begin{lstlisting}[style=coloredverbatimError]
// the element of a generator comprehension cannot yield, it is the value the generator yields
fn main () {
    let _ = (yield { yield x; x * 2 } for x in 0 .. 3);
}
\end{lstlisting}

\begin{lstlisting}[style=bashVerb, breaklines=true, escapechar=~]
Error[E4301] : cannot yield out of the element of a generator comprehension
 --> test_resources/generators/test68.yr:(3,22)
 3  ~\lstbox{\textSFxi{}}~     let _ = (yield { yield x; x * 2 } for x in 0 .. 3);
    ~\lstbox{\textSFv{}}~                      ^^^^^
    ~\lstbox{\textSFxi{}}~
    = when validating test68::main -- (test_resources/generators/test68.yr:(2,4))
    ~\lstbox{\textSFii{}}~~\lstbox{\textSFx{}}~~\lstbox{\textSFx{}}~~\lstbox{\textSFx{}}~~\lstbox{\textSFx{}}~~\lstbox{\textSFx{}}~~\lstbox{\textSFx{}}~~\lstbox{\textSFx{}}~
\end{lstlisting}

\subsubsection*{How to fix it}

Test the condition outside the comprehension, or write the generator as a function, \texttt{yield f (it: …)-\textgreater{} T \{ for x in it \{ … yield e; \} \}}, where a loop, a \tokennolst{continue}, a \tokennolst{return} and a \tokennolst{yield} have their usual meaning. That is also the way to yield several values per element.

\errorcode{E4302}{the size of a list comprehension filtered by if is not determinable at compile time}
\label{sec:codes:e4302}

Raised during \textbf{validation}, from \texttt{ValidateErrorMessage::\allowbreak{}LIST\_\allowbreak{}COMPR\_\allowbreak{}FILTER\_\allowbreak{}SIZE} (\textit{src/\allowbreak{}ymirc/\allowbreak{}semantic/\allowbreak{}validator/\allowbreak{}errors.yr}).

\subsubsection*{What it means}

A comprehension filtered by \tokennolst{if} collects the values the filter keeps, and which ones those are is only known once the iteration ran. Its result is therefore a slice, appended to as the iteration goes, and a slice is allocated on the heap: the comprehension has to be enclosed in a \tokennolst{copy}.

Without the filter the same comprehension has the length of the iteration, so it can be an array on the stack, which is why dropping the \tokennolst{copy} is accepted there and rejected here.

A comprehension over a tuple is the exception: its iteration is unrolled during validation, so the filter is computed there and the length of the result is known. It stays an array on the stack, and needs no \tokennolst{copy} -- at the price of a filter that has to be a compile time value, or \hyperref[sec:codes:e4303]{E4303} is raised.

This is about the \textit{number} of values, not about their type: a filter whose test is not a boolean raises \hyperref[sec:codes:e4095]{E4095} instead, and a comprehension whose iteration itself has no compile time length raises \hyperref[sec:codes:e4109]{E4109}.

\subsubsection*{Example}

\textit{test\_\allowbreak{}resources/\allowbreak{}diagnostics/\allowbreak{}test146.yr}:

%% check: skip
\begin{lstlisting}[style=coloredverbatimError]
fn main () {
    let _ = [i for i in 0 .. 10 if i != 6];
}
\end{lstlisting}

\begin{lstlisting}[style=bashVerb, breaklines=true, escapechar=~]
Error[E4302] : the size of a list comprehension filtered by if is not determinable at compile time
 --> test_resources/diagnostics/test146.yr:(2,38)
 2  ~\lstbox{\textSFxi{}}~     let _ = [i for i in 0 .. 10 if i != 6];
    ~\lstbox{\textSFv{}}~                                      ^^
    ~\lstbox{\textSFxi{}}~
    = help: did you mean to enclose it within a copy?
    ~\lstbox{\textSFxi{}}~
 2  -     let _ = [i for i in 0 .. 10 if i != 6];
 2  +     let _ = copy [i for i in 0 .. 10 if i != 6];
    ~\lstbox{\textSFxi{}}~
    = when validating test146::main -- (test_resources/diagnostics/test146.yr:(1,4))
    ~\lstbox{\textSFii{}}~~\lstbox{\textSFx{}}~~\lstbox{\textSFx{}}~~\lstbox{\textSFx{}}~~\lstbox{\textSFx{}}~~\lstbox{\textSFx{}}~~\lstbox{\textSFx{}}~~\lstbox{\textSFx{}}~
\end{lstlisting}

\subsubsection*{How to fix it}

Enclose the comprehension in a \tokennolst{copy}, as the fix of the diagnostic suggests:

%% check: skip
\begin{lstlisting}[style=coloredverbatim]
let _ = copy [i for i in 0 .. 10 if i != 6];
\end{lstlisting}

\errorcode{E4303}{the filter of the list comprehension, constructed from type \errhole{}, must be determinable at compile time}
\label{sec:codes:e4303}

Raised during \textbf{validation}, from \texttt{ValidateErrorMessage::\allowbreak{}LIST\_\allowbreak{}COMPR\_\allowbreak{}FILTER\_\allowbreak{}CTE} (\textit{src/\allowbreak{}ymirc/\allowbreak{}semantic/\allowbreak{}validator/\allowbreak{}errors.yr}).

\subsubsection*{What it means}

A comprehension written between parentheses constructs a tuple, and a tuple has an arity fixed at compile time. Its iteration is unrolled during validation, so the filter choosing which elements the tuple holds is computed there too, and a test that is not a compile time value leaves the arity unknown.

What a filter can test therefore depends on what the iteration makes available at compile time:

\begin{itemize}
\item over a range, the value of the iteration is a literal, so it can be tested,
\item over a tuple, both the index and the element are, so both can be tested,
\item over a slice or an array, only the index is: the elements live in memory the comprehension does not read at compile time.
\end{itemize}

A comprehension written between brackets escapes the constraint only when its iteration is not unrolled: it then builds a slice at run time, where the filter is an ordinary test -- see \hyperref[sec:codes:e4302]{E4302}. Over a tuple it builds an array whose length is fixed at compile time, so the filter is computed at compile time there as well, and this error is raised on that form too.

\subsubsection*{Example}

\textit{test\_\allowbreak{}resources/\allowbreak{}diagnostics/\allowbreak{}test147.yr}:

%% check: skip
\begin{lstlisting}[style=coloredverbatimError]
fn main (args: [[c8]]) {
    let n = args.len;
    let _ = (i for i in 0 .. 5 if i != n,);
}
\end{lstlisting}

\begin{lstlisting}[style=bashVerb, breaklines=true, escapechar=~]
Error[E4303] : the filter of the list comprehension, constructed from type (..i32), must be determinable at compile time
 --> test_resources/diagnostics/test147.yr:(3,37)
 3  ~\lstbox{\textSFxi{}}~     let _ = (i for i in 0 .. 5 if i != n,);
    ~\lstbox{\textSFv{}}~                                     ^^  error: value of type bool is needed but unknown at compilation time
    ~\lstbox{\textSFxi{}}~
    = when validating test147::main -- (test_resources/diagnostics/test147.yr:(1,4))
    ~\lstbox{\textSFii{}}~~\lstbox{\textSFx{}}~~\lstbox{\textSFx{}}~~\lstbox{\textSFx{}}~~\lstbox{\textSFx{}}~~\lstbox{\textSFx{}}~~\lstbox{\textSFx{}}~~\lstbox{\textSFx{}}~
\end{lstlisting}

\subsubsection*{How to fix it}

Test something the comprehension knows at compile time, or iterate something whose length is not fixed at compile time, so that the result is a slice built at run time:

%% check: skip
\begin{lstlisting}[style=coloredverbatim]
let _ = copy [i for i in 0 .. 5 if i != n];
\end{lstlisting}
